\section*{Capítulo \ref{ch:g2:drawing-to-map} --- Do desenho ao mapa}

\begin{solution}{exo:g2:drawing-to-map:1}
Como vista de cima.
\end{solution}

\begin{solution}{exo:g2:drawing-to-map:2}
As linhas azuis na parede oeste são as janelas. No canto nordeste fica o
armário.
\end{solution}

\begin{solution}{exo:g2:drawing-to-map:3}
Um título, uma \omterm{def:g2:drawing-to-map:legend}{legenda}, uma \omterm{def:g2:drawing-to-map:scale}{escala} e uma \omterm{def:g2:drawing-to-map:directions}{rosa dos ventos} (ou uma seta que aponta para o norte).
\end{solution}

\begin{solution}{exo:g2:drawing-to-map:4}
O norte fica no alto; o oeste fica à esquerda.
\end{solution}

\begin{solution}{exo:g2:drawing-to-map:5}
Um quadradinho representa 1 metro: a carteira tem 2 metros de comprimento e 1 metro de largura.
\end{solution}

\begin{solution}{exo:g2:drawing-to-map:6}
O quadradinho marrom é o sítio; a mancha verde é o bosque.
\end{solution}

\begin{solution}{exo:g2:drawing-to-map:7}
O bosque fica ao sul do \omterm{def:g1:landscapes-around-us:water}{rio} e a leste do sítio.
\end{solution}

\begin{solution}{exo:g2:drawing-to-map:8}
Um quadradinho representa 1 quilômetro: 10 quilômetros de oeste a leste e 8
quilômetros de sul a norte.
\end{solution}

\begin{solution}{exo:g2:drawing-to-map:9}
Três fileiras de três carteiras: 9 carteiras. A porta fica na parte sul da
parede leste; então a carteira mais perto da porta fica no sul (no canto
sudeste da sala).
\end{solution}

\begin{solution}{exo:g2:drawing-to-map:10}
A cama tem 2 quadradinhos de comprimento. O quarto tem 4 quadradinhos por 3
quadradinhos: 3 fileiras de 4 quadradinhos, ou seja, $4 + 4 + 4 = 12$ quadradinhos.
\end{solution}

\begin{solution}{exo:g2:drawing-to-map:11}
``Esquerda'' depende de para onde Lena está virada: se ela der meia-volta,
o \omterm{def:g1:landscapes-around-us:water}{rio} fica à sua direita. Para todo mundo: ``o \omterm{def:g1:landscapes-around-us:water}{rio} fica ao norte do sítio''.
\end{solution}
